% Plantilla realizada por Santiago Morante Cendrero

%parametros de tipo libro
\documentclass[10pt,a4paper]{book}

%idioma espa�ol y acentos
\usepackage[spanish]{babel}
\usepackage[latin1]{inputenc}

%algunos s�mbolos matem�ticos y paquete para usar subim�genes
\usepackage{amsmath}
\usepackage{amsfonts}
\usepackage{amssymb}
\usepackage{graphicx}
\usepackage{subfigure}
\usepackage{listings}
\usepackage{appendix}

%para generar �ndice con hiperv�nculos
\usepackage{hyperref}

%parametros del documento (sus propiedades)
\hypersetup{
    pdftitle={Arturo Arranz Mateo - TFG - a�o},
    pdfsubject={TFG - a�o},
    pdfauthor={Nombre del alumno},
    pdfkeywords={palabraclave1} {palabraclave2} {palabraclave3},
    colorlinks,
    citecolor=black,
    filecolor=black,
    linkcolor=black,
    urlcolor=black,
}

%empieza el documento
\begin{document}  

%elementos antes del trabajo en s� se meten dentro de frontmatter
\frontmatter

%cada incluye referencia a un archivo de tipo .tex
\begin{titlepage}
\begin{center}

%forma de introducir im�genes. el \\[0.5 cm] de final de l�nea introduce un salto de ese tama�o.
%width=1\textwidth indica el tama�o de la im�gen (valores entre 0-1). 
 \includegraphics[width=1\textwidth]{figuras/cabecera.png}  \\[0.5 cm]

\large \textsc{Departamento de Ingenier�a Inform�tica} \\ [1 cm]

\large TRABAJO FIN DE GRADO\\[1 cm]

\huge \textsc{Reconocimento de acciones mediante kinect y aprendizaje autom�tico}\\[7 cm]

%flushleft alinea a la izquierda el texto
\begin{flushleft} \Large
\emph{Autor:} Arturo Arranz Mateo\\[0.5 cm]
\emph{Director:} nombre del director \\
\emph{Tutor:} nombre del tutor
\end{flushleft}

%rellena de blanco el resto de la p�gina para escribir abajo del todo
\vfill

% Bottom of the page
{\large Ciudad, Mes a�o}

\end{center}
\end{titlepage}

\begin{flushleft}

Copyright \copyright  a�o. Nombre del alumno

%ejemplo de licencia, se puede elegir cualquier otra

Esta obra est� licenciada bajo la licencia Creative Commons Atribuci�n-NoComercial-SinDerivadas 3.0 Unported (CC BY-NC-ND 3.0). Para ver una copia de esta licencia, visite http://creativecommons.org/licenses/by-nc-nd/3.0/deed.es o env�e una carta a Creative Commons, 444 Castro Street, Suite 900, Mountain View, California, 94041, EE.UU.

Todas las opiniones aqu� expresadas son del autor, y no reflejan necesariamente las opiniones
de la Universidad Carlos III de Madrid.

\end{flushleft}

\cleardoublepage

\begin{flushleft} \large
\textbf{T�tulo:} t�tulo del trabajo \\
\textbf{Autor:} nombre del alumno\\
\textbf{Director:} nombre del director \\ 
\textbf{Tutor:} nombre del tutor\\ [1 cm]

\end{flushleft} 

\begin{center} \LARGE
EL TRIBUNAL \\ [1 cm]
\end{center}

\begin{flushleft} \LARGE
Presidente: \\ [1 cm]
Vocal: \\ [1 cm]
Secretario: \\ [1.5 cm]
\end{flushleft}

\large
Realizado el acto de defensa y lectura del Trabajo Fin de Grado el d�a ....... de ....................   de ... en .........., en la Escuela Polit�cnica Superior de la Universidad Carlos III de Madrid, acuerda otorgarle la CALIFICACI�N de: \\ [2 cm]

\begin{center}
 \large VOCAL \\ [2.2 cm]
\end{center}

\begin{minipage}{0.5\textwidth}
 \begin{flushleft}
 \large SECRETARIO
\end{flushleft}
\end{minipage}
\begin{minipage}{0.5\textwidth}
\begin{flushright}
 \large PRESIDENTE
\end{flushright} 
\end{minipage}

\chapter{Agradecimientos}

Agradezco a ............

%chapter introduce un nuevo cap�tulo
\chapter{Resumen}

Este proyecto de resume en.................

\paragraph{Palabras clave:} palabraclave1, palabraclave2, palabraclave3.

\chapter{Abstract}

In this project...

\paragraph{Keywords:} keyword1, keyword2, keyword3.

%genera �ndice
\tableofcontents

%�ndice de figuras. Tambi�n se podr�a hacer uno de tablas (listoftables)
\listoffigures
 
%empieza la parte descriptiva del trabajo
\mainmatter
 
\chapter{Introducci�n}

En este cap�tulo...

%section es un apartado dentro de un chapter. Tambi�n existe subsection y subsubsection
\section{Estado del arte}

%las referencias a art�culos se ponen con \cite, 
%las referencias a im�genes \ref, 
%y las referencias a ecuaciones \eqref

Este tema.... Esto es un ejemplo de cita de un art�culo \cite{gonzalezeducational}.

%itemize es una lista. Cada t�rmino lleva delante un \item
\begin{itemize}
\item \textbf{ejemplo de lista de puntos}. Ejemplo.
\item \textbf{ejemplo2 de lista}. Ejemplo2.
\end{itemize} 

Ejemplo de referencia a figura (figura \ref{uc3m}).

%caption es el pie de foto, y label es el nombre que se da a la imagen para referenciarla despu�s. label no puede llevar acentos y no se muestra de cara al documento final (es s�lo interno).
\begin{figure}[h]
\centering
\includegraphics[width=0.45\textwidth]{figuras/uc3m}   
\caption{Logotipo de la UC3M \copyright UC3M}
\label{uc3m}
\end{figure}


\section{Motivaci�n del proyecto}

La idea...

\section{Estructura del documento}

A continuaci�n y para facilitar la lectura del documento, se detalla el contenido de cada cap�tulo.

\begin{itemize}
\item En el cap�tulo 1 se realiza una introducci�n.
\item En el cap�tulo 2 se hace un repaso...
\end{itemize}

  
  %partes finales del trabajo: conclusiones, bibliografia y anexos
\backmatter

\chapter{Conclusiones}

Se presentan a continuaci�n las conclusiones...

\section{Conclusi�n}

Una vez finalizado el proyecto...

\section{Desarrollos futuros}

Un posible desarrollo...


\chapter{State of art}
\\
ARtajsdasdnasdfjaskodfcas
sdgfasdf

sfgad
fhsf
hgs
ghs
fghs
gs
afg
dfg
sdfg
dfg

%estilo de bibliograf�a: plana, alfa...
\bibliographystyle{plain}

%genera doble hoja en blanco
\cleardoublepage

%apartado de bibliograf�a
\addcontentsline{toc}{chapter}{Bibliograf�a}

%se incluye la bibliograf�a. Archivo de tipo .bib (bibtex)
\bibliography{bibliografia}

%fin del documento
\end{document}